\section{Introducción}
Los exoesqueletos de mano son dispositivos robóticos de asistencia y
rehabilitación cuya función es acompañar o restituir el movimiento de los dedos.
Para que la asistencia sea cómoda y segura, el mecanismo debe reproducir con
fidelidad las trayectorias articulares naturales del dedo a lo largo de todo el
rango de flexión. Lograrlo depende en gran medida de elegir correctamente las
dimensiones de los eslabones del mecanismo, un problema conocido como
\emph{síntesis dimensional}.

La síntesis dimensional de mecanismos para generación de trayectoria es un
problema de optimización no lineal, multimodal y fuertemente restringido: el
espacio de soluciones contiene amplias regiones donde el mecanismo
simplemente no ensambla o invierte su movimiento, y la relación entre las
longitudes y la curva resultante es altamente no convexa. Estas
características dificultan el uso de métodos clásicos basados en gradiente y
motivan el empleo de metaheurísticas poblacionales como la Evolución
Diferencial (ED) y los Algoritmos Genéticos (AG).

Aunque ambas familias se usan ampliamente, su desempeño relativo depende de la
naturaleza del problema, de la codificación y del ajuste de hiperparámetros. La
elección entre una u otra rara vez es obvia y suele requerir evidencia empírica
sobre la instancia concreta. En este trabajo abordamos esa pregunta para un
mecanismo de exoesqueleto de dedo real, formulando una comparación controlada
en la que ambos algoritmos comparten exactamente el mismo planteamiento,
función objetivo y presupuesto de generaciones.

Para que la comparación sea robusta y reproducible, se optimiza una versión
\emph{parcial} del mecanismo, que reproduce el movimiento hasta la falange
medial (articulación IFD). Esta instancia es más pequeña que el modelo completo
de tres falanges, pero conserva las dificultades esenciales del problema
(acoplamiento de etapas, restricciones de ensamble y monotonicidad), lo que la
hace idónea para aislar el efecto del algoritmo.

\subsection{Contribuciones}
Las contribuciones de este artículo son: (i) una formulación unificada del
problema de síntesis dimensional parcial con 16 parámetros, función objetivo
basada en distancia de Chamfer ponderada y penalizaciones de viabilidad;
(ii) una comparación controlada de ED y AG bajo condiciones idénticas, evaluando
no sólo el error de ajuste sino también la coherencia y compacidad del
mecanismo resultante; y (iii) un análisis de las razones por las que la ED
supera al AG en este dominio continuo. El modelo cinemático y el diseño del
mecanismo empleados como base fueron desarrollados por Pavel Palma Nicolás, ex
alumno de la universidad de los autores.
